\section{Training}
\label{sec:training}
\subsection{Pre-training Recipe}

\subsubsection{Dual-Stream Co-Training}
We train \ours on two complementary data streams simultaneously. The VLA stream is built from the full multi-source manipulation corpus described in \cref{sec:data}, comprising real-robot demonstrations, egocentric human-hand manipulation videos, and human-to-robot synthesized trajectories. 

The VLM stream consists of large-scale vision-language supervision data (\cref{sec:vlm_data}), co-trained alongside the VLA stream to prevent the pre-trained perceptual and language capabilities from degrading under action prediction optimization, which would directly weaken the model's ability to interpret novel instructions and generalize to unseen visual contexts~\citep{driess2025knowledge}. In practice, we adopt a 9:1 ratio of robot data to VL data.

% We train \ours on two complementary data streams simultaneously: a VLA stream for action learning and a VLM stream for preserving and extending general vision-language capabilities.

% The VLA stream is built from the full multi-source manipulation corpus described in \cref{sec:data}. It includes real-robot demonstrations collected from diverse robot embodiments, egocentric human-hand manipulation videos, and human-to-robot synthesized trajectories. Each VLA sample contains visual observations, proprioceptive states, action chunks, and a natural-language task instruction.

% The VLM stream consists of large-scale vision-language supervision data \cref{sec:vlm_data}. We co-train on this stream together with the VLA data so that the model retains its pretrained visual and language capabilities during VLA training~\citep{driess2025knowledge}. Optimization on action prediction alone can degrade general-purpose perception and language understanding, which in turn weakens the VLA model's ability to interpret novel instructions and generalize to unseen visual contexts. In practice, we adopt a 9:1 ratio of robot data to VL data for co-training.

\subsubsection{Training Objectives}

\textbf{Flow matching loss.}
For each VLA sample, we are given a ground-truth action chunk $\mathbf{a} \in \mathbb{R}^{T \times D}$. Following the flow-matching formulation, we construct a noisy interpolant
\[
\mathbf{x}_t = (1-t)\,\boldsymbol{\epsilon} + t\,\mathbf{a},
\]
where $\boldsymbol{\epsilon} \sim \mathcal{N}(\mathbf{0}, \mathbf{I})$ and $t \sim \mathrm{Beta}(1, 1.5)$. The action expert is trained to predict the corresponding velocity field $\mathbf{v} = \mathbf{a} - \boldsymbol{\epsilon}$, minimizing:
\begin{equation}
    \mathcal{L}_{\mathrm{FM}} = \mathbb{E}_{\mathbf{a},\,\boldsymbol{\epsilon},\,t}
        \left\| f_\theta\!\left(\mathbf{x}_t,\, t,\, \mathbf{s},\, \mathbf{o}\right) - \left(\mathbf{a} - \boldsymbol{\epsilon}\right) \right\|_2^2,
    \label{eq:fm_loss}
\end{equation}
where $f_\theta$ denotes the full model conditioned on the proprioceptive state $\mathbf{s}$ and the visual-language observation $\mathbf{o}$. Gradients from $\mathcal{L}_{\mathrm{FM}}$ are applied to both the VLM backbone and the action expert.

Because different robot embodiments populate different subsets of the 80-dimensional canonical action space (\cref{sec:unified_state_action}), 
we apply a composed binary mask $\mathbf{m} \in \{0,1\}^{T \times D}$ that restricts the objective to only the dimensions and time steps carrying valid supervision.
The mask is constructed from three complementary sources.
The \emph{per-dimension slot mask} identifies actively populated dimensions for the current embodiment. For instance, a single-arm gripper populates the joint, end-effector, and gripper fields of one arm while leaving the opposite arm and hand slots as zeros.
The \emph{step validity mask} excludes time steps flagged as anomalous by the data curation pipeline (\cref{sec:state_action_filtering}) or outside episode boundaries, with all subsequent steps also masked once any step is deemed invalid to preserve causal consistency.
For egocentric human data, a \emph{per-hand validity mask} zeros out an entire arm slot from the moment the corresponding hand exits the camera view, preventing the model from being trained on occluded hand trajectories.
The three masks are AND-combined, and the masked flow matching loss becomes a per-sample average over valid entries only:
\begin{equation}
    \mathcal{L}_{\mathrm{FM}} = \frac{1}{B}\sum_{i=1}^{B}
        \frac{\sum_{t,j}\, m_{i,t,j}\,\bigl(f_\theta(\mathbf{x}_{i,t},\, t_i,\, \mathbf{s}_i,\, \mathbf{o}_i)_{j} - v_{i,t,j}\bigr)^2}
             {\sum_{t,j}\, m_{i,t,j}},
    \label{eq:fm_loss_masked}
\end{equation}
where $B$ is the batch size and subscripts $t,j$ index the time step and dimension.
This formulation ensures that every sample in the batch contributes equally to the gradient regardless of how many dimensions are active, preventing embodiments with more populated slots from dominating optimization.

\textbf{VLM next-token prediction loss.}
For each VLM sample, the backbone is trained with the standard autoregressive next-token prediction objective:
\begin{equation}
    \mathcal{L}_{\mathrm{VLM}} = -\mathbb{E} \sum_{i} \log p_\phi\!\left(y_i \mid y_{<i},\, \mathbf{c}\right),
    \label{eq:vlm_loss}
\end{equation}
where $y_i$ is the target response token at position $i$, and $\mathbf{c}$ is the input context, which may be text-only or an interleaved sequence of images and text.
The overall objective is
\begin{equation}
    \mathcal{L} = \mathcal{L}_{\mathrm{FM}} + \lambda\,\mathcal{L}_{\mathrm{VLM}},
    \label{eq:total_loss}
\end{equation}
where $\lambda$ controls the relative weight of the two losses.
We set $\lambda = 0.1$ so that VLM supervision provides stabilizing regularization without overshadowing action learning. Separate learning rates are used for the backbone and the action expert to account for their different initialization scales.
To amortize the cost of the VLM forward pass, the action expert performs $K_{\mathrm{repeat}}{=}8$ repeated diffusion steps per training sample, drawing 8 independent noise samples and timesteps for the same action chunk, and substantially improving training efficiency without increasing data consumption.


\subsection{Post-Training Recipe}

\subsubsection{Domain-specific Supervised Fine-tuning}
\label{sec:domain_sft}
Once the foundation model has been pre-trained on the full heterogeneous corpus, we adapt it to specific deployment scenarios through supervised fine-tuning (SFT). 
Rather than training specialist policies for individual tasks, we adopt a \emph{generalist SFT} paradigm. 
For each benchmark or real-world deployment scenario, all available demonstration data is combined into a single training set, 
producing one unified fine-tuned model that can execute every task within the target domain. 
%This avoids the overhead of maintaining multiple specialist checkpoints and, as we show in \cref{sec:experiment}, 
%yields models whose cross-task performance matches or exceeds that of per-task specialists.

Compared with pre-training, SFT differs in several aspects.
First, the SFT procedure only optimizes the flow matching objective $\mathcal{L}_{\mathrm{FM}}$ of \cref{eq:fm_loss_masked} without the VLM next-token prediction loss.
We disable the multi-stage data curation filtering of \cref{sec:state_action_filtering} and train on the complete unfiltered data, preserving every valid demonstration for post-training.
We apply color jitter augmentation to the input images.
In addition, with the model initialized from a pre-trained checkpoint, SFT is conducted on fewer GPUs with fewer training steps than pre-training. 

\subsubsection{Co-Training in Post-Training}
Domain-specific SFT has become an important protocol for quantitatively evaluating the quality of a pre-trained VLA model. 
A strong pre-trained model is expected to adapt efficiently to a target domain after fine-tuning on the benchmark training set. However, this evaluation protocol can also expose a critical failure mode. After extensive benchmark-specific SFT, a VLA model may improve task performance by exploiting repeated visual and task patterns in the benchmark, while becoming less sensitive to the language instruction. In this case, the policy is no longer strongly conditioned on both vision and language. Instead, it behaves more like a vision-action pattern matcher. We refer to this phenomenon as VLA-to-VA degradation.

This degradation stems from three compounding factors. First, domain-specific SFT datasets are limited in diversity with concentrated visual layouts and instruction expressions, making shortcut correlations between scene patterns and actions easy to learn. Second, training and test splits share similar visual patterns, allowing high benchmark scores to be achieved through pattern memorization rather than genuine language grounding. As a result, such benchmarks overestimate instruction-following ability and fail to distinguish true language-conditioned control from benchmark-pattern overfitting. Third, models without strong compositional grounding ability tend to treat language as a weak context signal during SFT, with actions increasingly dominated by visual shortcuts learned from the benchmark data.

To mitigate this risk, we propose a mixed post-training strategy as an optional enhancement to the standard domain SFT described in \cref{sec:domain_sft}: rather than fine-tuning solely on the benchmark training set, the model is co-trained with a subset of pre-training data filtered by distributional proximity to the target domain. 
This provides broader adaptation signals while preserving the base model's robust execution capabilities, without introducing unrelated data that might dilute domain-specific learning. 
In our main experiments, we follow the standard domain-specific SFT protocol to ensure fair comparison with baselines; the mixed post-training strategy is validated as an additional enhancement in \cref{sec:mixposttrain}.

To directly evaluate whether the language-following capability is preserved, we construct a new benchmark, RoboTwin-IF (\textbf{I}nstruction \textbf{F}ollowing), based on RoboTwin 2.0~\citep{chen2025robotwin}. 
The benchmark examines whether the model performs the instructed action in the same or similar visual scenes with different instructions, rather than relying on visual pattern matching to select a default behavior. 
The detailed benchmark design is presented in \cref{sec:eval_protocol}.

% To mitigate this, rather than fine-tuning on the benchmark training set alone, we co-train with both benchmark data and a selected subset of pretraining data chosen for its proximity to the target benchmark distribution. This is not intended to introduce a broad and unrelated data mixture, but provides additional adaptation signals beyond the narrow official training split while preserving the robust physical execution capability of the foundation model.

% To better evaluate language conditioning, we construct a new benchmark, RoboTwin-IF (\textbf{I}nstruction \textbf{F}ollowing) , based on RoboTwin 2.0 ~\citep{chen2025robotwin}. Rather than only testing task completion, the benchmark examines whether the model performs the instructed action and satisfies the corresponding success criteria in the same or similar scenes with different instructions. Specifically, we implement task variations along three dimensions: action type, manipulated object, and target location. By recombining these factors into diverse instructions, the benchmark is designed to measure whether the policy genuinely uses language as a control signal rather than relying on visual pattern matching. More details and results are present in Section~\ref{sec:mixposttrain}.

% As a result, such benchmarks may overestimate instruction-following ability and fail to distinguish true language-conditioned control from benchmark-pattern overfitting. Third, this issue is further amplified when the pre-trained model lacks strong language-conditioned policy modulation and compositional grounding ability. Without these capabilities, language may become a weak contextual signal during SFT, while the final action is dominated by visual shortcuts learned from the benchmark data.

% To mitigate this problem, we develop a mixed post training strategy on top of the strong Qwen Manipulation foundation model. Instead of adapting the model only on the benchmark provided training set, we cotrain it with both benchmark data and a selected subset of pretraining data. The additional data is selected according to its proximity to the target benchmark distribution. It is not intended to introduce a broad and unrelated data mixture, but to provide benchmark relevant post training signals beyond the narrow official training split. Built upon the strong Qwen Manipulation foundation model, this training procedure allows the model to absorb these additional adaptation signals while preserving robust physical execution capability. As a result, the model can improve its target domain performance without collapsing into benchmark specific visual pattern matching.


% To better evaluate this capability, we construct an instruction-following benchmark that explicitly separates task completion from instruction-conditioned task completion. Rather than only testing whether the final task is completed, the benchmark examines whether the model performs the action specified by the instruction under varied task, object, target, and language conditions. By recombining semantic factors into diverse instructions and reducing reliance on repeated train-test patterns, the benchmark is designed to measure whether the policy truly uses language as a control signal, instead of solving the task through visual pattern matching. We present the empirical analysis in Section X.
% While specialist SFT improves performance on target-domain tasks, it often comes at the cost of reduced generalization compared with generalist SFT. In particular, fine-tuning a pre-trained VLA model on a small amount of domain-specific data can lead to a noticeable degradation in instruction-following ability. Although the model adapts quickly to the target domain, the limited diversity of domain data—especially in instruction expressions and constraint combinations—makes it prone to overfitting. As a result, the model may learn to complete the task itself while becoming less sensitive to explicit linguistic constraints in the instruction. This issue is especially critical in embodied settings, where correct execution depends not only on task completion, but also on whether the action is carried out by the correct agent, with the correct object, and toward the correct goal.

% To address this issue, we adopt a joint training strategy that mixes domain-specific data with a subset of general pre-training data during adaptation. Instead of fine-tuning only on the target domain, the model is trained on both data sources in a unified framework. The domain data provides task- and environment-specific supervision, while the pre-training data preserves broader instruction diversity, compositional coverage, and language-conditioned behavior patterns. In our setting, the model is initialized from a pre-trained backbone, and the adaptation stage is performed by co-training on the two data distributions. This mixed training procedure improves domain specialization while mitigating the loss of general instruction-following capability.

% To better quantify this effect, we construct a benchmark specifically designed to measure instruction-following behavior in embodied tasks. The benchmark is built by decomposing each task into several key semantic factors, including the acting agent, action type, manipulated object, and target location or target state. We then generate randomized instructions by recombining these factors into diverse language conditions. This design allows us to systematically test whether the model follows explicit constraints in the instruction, rather than merely achieving the final task outcome. In other words, the benchmark distinguishes task success from instruction-conditioned task success.

% In a dual-arm manipulation setting, a model fine-tuned only on limited domain data is often able to successfully grasp the target object, but may ignore the instruction constraint specifying which arm to use. For example, when the instruction explicitly requires the left arm, the fine-tuned model may still execute the action with the right arm. After introducing mixed pre-training data in the joint training stage, this behavior is substantially improved: the model becomes significantly more consistent in following explicit constraints such as left-versus-right arm selection. This result suggests that joint training is an effective way to mitigate instruction-following degradation caused by small-scale domain adaptation.
